\documentclass[11pt]{article}
\usepackage{amsmath,amssymb,amsthm}
\usepackage[margin=1.1in]{geometry}
\newtheorem{theorem}{Theorem}
\newtheorem{lemma}{Lemma}
\newtheorem{definition}{Definition}
\newtheorem{remark}{Remark}
\newcommand{\R}{\mathbb{R}}
\title{Dead neurons resurrect by an entry--exit law:\\ the smoothed-ReLU limit of
training dynamics at an activation boundary\\ \large DRAFT v1 --- 2026-07-24 ---
not for circulation; pending adversarial review}
\author{L.~Caloi Santos}
\date{}
\begin{document}
\maketitle

\begin{abstract}
For shallow softplus networks under gradient flow we show that, in the sharp
limit $\beta\to\infty$ of the activation stiffness, the death and resurrection
of a neuron is governed by a hybrid system \emph{with memory}: the dead
neuron's parameters freeze, and reactivation occurs when the signed time
integral of the network residual returns to zero --- the classical entry--exit
function of slow--fast dynamics --- with rate $O(\beta^{-1})$. The switching
face carries an identically zero field on the dead side, a degeneracy outside
generic piecewise-smooth exit theory and precisely of the exponentially flat
type of Kristiansen; the coupling of the forcing to the flow itself
(the residual is generated by the very trajectory it gates) has, to our
knowledge, no analogue in the entry--exit literature. The law is independent
of the regularization within monotone sigmoid classes, including a GELU-type erf smoothing. All constants are verified numerically;
companion pre-registered experiments probe the stochastic gate under SGLD.
\end{abstract}

\section{Setup and hypotheses}
Data $(x_i,y_i)_{i=1}^n$ in general position, shallow network
$f_\theta(x)=\sum_k g_\beta(w_kx+b_k)$ with $g_\beta(z)=\beta^{-1}\log(1+e^{\beta z})$
(softplus; see \S\ref{sec:reg} for general sigmoids), quadratic loss
$L=\tfrac1{2n}\sum_i(f_\theta(x_i)-y_i)^2$, gradient flow $\dot\theta=-\nabla L$.
Fix a neuron $A$ and a datum $j$; write $z=\theta_A\cdot(x_j,1)$ for its margin,
$c=\|(x_j,1)\|^2/n$, and $R(t)=y_j-f_{-A}(x_j;\theta_B(t))$ for the residual
\emph{excluding} $A$. The \emph{hybrid reference} $\theta^{hyb}$ runs the full
ReLU flow until the first time $t_{hit}$ with $z=0$, then freezes $\theta_A$ and
runs $\theta_B$ alone; set $R^{hyb}$, $A^{hyb}(t)=\int_{t_{hit}}^t R^{hyb}$.

\smallskip\noindent\textbf{Hypotheses.}
\textbf{(H1)} $R^{hyb}(t_{hit})<0$ (transversal death; this also gives the
transversal arrival, Lemma~\ref{lem:corner-in}).
\textbf{(H2)} $A^{hyb}$ has a first zero at $t_{exit}<T$ with
$R^{hyb}(t_{exit})>0$ and $A^{hyb}<0$ on $(t_{hit},t_{exit})$.
\textbf{(H3)} the other margins of $A$ stay $\le-\delta_1<0$ on the episode.
\textbf{(H4)} no other switching event occurs in $[0,T]$: the margins of the
remaining neurons stay bounded away from zero along the hybrid reference, the
activation pattern of $A$ is constant on $[0,t_{hit})$, and no second episode
of $A$ begins before $T$.

\section{Main result}
\begin{definition}[hybrid limit]\label{def:hybrid}
$\theta^\infty$ equals: the ReLU flow on $[0,t_{hit}]$; on $[t_{hit},t_{exit}]$,
$\theta_A\equiv\theta_A(t_{hit})$ (frozen) and $\theta_B$ the $B$-only flow; on
$[t_{exit},T]$, the ReLU flow restarted at
$(\theta_A(t_{hit}),\theta_B^{hyb}(t_{exit}))$.
\end{definition}

\begin{theorem}[T2$'$-min]\label{thm:main}
Under (H1)--(H4), for every $\delta>0$,
\[
\sup_{t\in[0,T]\setminus(B_\delta(t_{hit})\cup B_\delta(t_{exit}))}
\|\theta^\beta(t)-\theta^\infty(t)\| = O_\delta\!\left(\tfrac{\log\beta}{\beta}\right),
\qquad
t^\beta_{release}=t_{exit}+O(\beta^{-1}),
\]
where $t^\beta_{release}$ is the first return of the margin $z$ to $0$ after
death. The $\log\beta$ is sharp and is carried solely by the normal coordinate
of $\theta_A$ inside the episode (the sticking depth of
Lemma~\ref{lem:release}); $\theta_B$, the tangential components of $\theta_A$,
and $t^\beta_{release}$ obey the cleaner rate $O(\beta^{-1})$.
The limit is that of the regularized flow: discrete ReLU gradient descent
freezes the dead neuron forever (the limits do not commute).
\end{theorem}

\section{Proof outline}\label{sec:proof}
Six lemmas; each has an executable verification script in the repository.

\begin{lemma}[normal form]\label{lem:gauge}
In the gauge $w=e^{\beta z}$, fast time $\tau=\beta t$:
$dw/d\tau=cR(t)\,w^2/(1+w)\,(1+E_\beta)$, $|E_\beta|\le\log(1+w)/(\beta|R|)$
--- the quadratic normal form of De Maesschalck--Schecter with
$\varepsilon=\beta^{-1}$. With $u=1/w$ and $\Phi(u)=\log u+u$ one has the exact
quadrature $d\Phi/dt=-c\beta R(t)(1+E_\beta)$.
\end{lemma}

\begin{lemma}[B-subsystem and feedback]\label{lem:gronwall}
On any episode interval where $w\le w_{\max}$, the coupling of $A$ into the
$B$-field is $\le\beta^{-1}\log(1+w_{\max})$ and
$\sup\|\theta_B^\beta-\theta_B^{hyb}\|\le C\beta^{-1}$ (Gr\"onwall). In the bulk
of sticking $u=O(\beta)$, so the coupling there is $O(\beta^{-2})$.
\end{lemma}

\begin{lemma}[entry corner]\label{lem:corner-in}
$\dot z^0(t_{hit}^-)=cR^{hyb}(t_{hit})<0$ by (H1); by comparison in the window
$z\in[\beta^{-1}\log w_0,\delta_z]$, the crossing of $\{w=w_0\}$ occurs at
$t_{hit}+O(\beta^{-1})$ with state $O(\beta^{-1})$-close to
$\theta^\infty(t_{hit})$.
\end{lemma}

\begin{lemma}[sticking and release, prescribed $R$]\label{lem:release}
For prescribed $R$ satisfying (H1)--(H2), with $A(t)=\int_{t_{hit}}^t R$: the
exact additive identity $d\Phi/dt=-c\beta R+c\log(1+w)$ holds (no division by
$R$; the interior zero of $R$ forced by (H1)--(H2) via the intermediate value
theorem is harmless --- the layer is held by the \emph{area}, not the sign of
$R$), and on compact subintervals of $(t_{hit},t_{exit})$,
$\Phi(u(t))=-c\beta A(t)(1+o(1))$, hence $u(t)=c\beta|A(t)|(1+o(1))$ (depth $z\sim-\beta^{-1}\log\beta$), and $u$
returns to $O(1)$ at $t_{exit}+O(\beta^{-1})$ --- no $\log\beta$ factor. Sub/super
solutions from the monotonicity of the flow in $R$ give explicit constants;
perturbations $\delta$ of $R$ shift $t_{exit}$ by
$\le c\,\delta\,(t_{exit}-t_{hit})/|A'(t_{exit})|+O(\delta^2)$.
\end{lemma}

\begin{lemma}[coupled bootstrap]\label{lem:bootstrap}
The set of times up to which ``$\theta_B$ is $2M\beta^{-1}$-close to the hybrid
and the margin stays in the layer'' is open and closed in
$[t_{cross},t_{exit}-\epsilon]$: under the assumption, coupling is
$O(\beta^{-1})$ (corners) $/$ $O(\beta^{-2})$ (bulk), so
Lemma~\ref{lem:gronwall} returns the better constant $M\beta^{-1}$ and
Lemma~\ref{lem:release} with $\delta=O(\beta^{-1})$ keeps the layer. Hence the
closed-loop law $t^\beta_{release}=t_{exit}+O(\beta^{-1})$.
\end{lemma}

\begin{lemma}[exit corner; scalar reduction]\label{lem:corner-out}
Post-release the field is transversal outward ($cR^{hyb}(t_{exit})>0$, (H2));
gluing as in Lemma~\ref{lem:corner-in}. Under (H3) the other data contribute
$O(e^{-\beta\delta_1})$ and the face-$j$ field is rank one along $(x_j,1)$, so
the tangential components of $\theta_A$ freeze to exponential accuracy and the
episode reduces to the scalar problem.
\end{lemma}

\section{Numerical verification}\label{sec:num}
Instance $X=(0.5,1.5)$, $Y=(0.5,1.0)$, $\theta_A(0)=(-0.5,0.27)$,
$\theta_B(0)=(1.144,0.261)$; $t_{hit}=0.1114$, $t_{exit}=1.1325$,
$\delta_1=0.508$. Selected checks (scripts \texttt{passo0}--\texttt{passo6}):
entry $t_{cross}=0.11141+7.13/\beta$ (fit residual $0.2\%$ over
$\beta\in[3\cdot10^2,10^4]$); depth ratio $u_{\max}/(c\beta|A|_{\max})$ increasing $0.888\to0.982$ toward
the predicted limit $1$ (prescribed-$R$ variant of the instance, whose own
$t_{exit}=1.1352$);
bulk feedback $\times\beta^2\approx43$ constant over $1.5$ decades; closed-loop
$\beta\,(t^\beta_{release}-t_{exit})\to-14.4$ (a definite first-order
coefficient, sharper than the theorem); tangential drift of $\theta_A$ at
machine zero.

\section{Regularization independence; discussion}\label{sec:reg}
The release criterion depends only on $\int R\,dt=0$, not on the sigmoid: for
the erf class ($g_\beta'=\Phi_{\rm gauss}(\beta z)$, a GELU-type smoothing;
GELU itself, $z\Phi_{\rm gauss}(z)$, has non-monotone derivative and is not a
member of the family for any $\beta$) the coupled system releases at the same $t_{exit}$ with
$O(\beta^{-1})$ gap (verified). Open: the cone apex (corank $\ge2$), where
weight decay produces genuine sliding toward the most singular stratum; the
SGLD gate (pre-registered experiments show the deterministic schedule is
falsified at first order by persistent residual fluctuations --- the stochastic
entry--exit is the object of the next cycle); and the non-commutation with discrete-time ReLU training (elementary here,
since $\mathrm{relu}'(0)=0$ freezes the dead neuron; the relation to the
broader non-uniqueness literature for ReLU gradient flow is deliberately
deferred to a verified revision). Conceptual anchors: entry--exit
function~\cite{DMS16,SDR23}; flat slow manifolds~\cite{Kri17}; generic exit
and regularized-limit theory~\cite{Jef16,GH}; nonsmooth training
formalism~\cite{BP21}; ReLU--LCP equivalence~\cite{APM20}; dying
ReLU~\cite{Lu19}; the canonical regularization of discontinuous fields is
Sotomayor--Teixeira~\cite{ST98,LST07} --- whose generic regions require both
one-sided normal components nonzero, so the dead face (one side identically
zero) is outside its hypotheses by construction.

\section*{Provenance}
Every lemma except the exit-corner half of Lemma~\ref{lem:corner-out} has a
dedicated executable symbolic/numeric verification committed alongside (the
exit corner is verified indirectly through the release times). Experiments were
adversarially reviewed (two rounds) and pre-registered by commit before
execution; the referee of this draft questioned the audit trail and the check
was re-run: the pre-registration commit precedes the confirmatory run in the
repository history. Two earlier experimental designs and one earlier conjecture
were rejected by the same process and are archived. All numerics are a single
two-neuron instance; breadth is future work.

\begin{thebibliography}{9}
\bibitem{DMS16} P.~De Maesschalck, S.~Schecter, \emph{The entry--exit function
and geometric singular perturbation theory}, J.~Diff.~Eq.\ 260 (2016) 6697--6715.
\bibitem{Kri17} K.~U. Kristiansen, \emph{Blowup for flat slow manifolds},
arXiv:1603.01821 (Nonlinearity, 2017).
\bibitem{GH} N.~Guglielmi, E.~Hairer, \emph{On the limit solution of regularized
piecewise-smooth dynamical systems}, preprint (Univ.\ Geneva); foundational
results in SIAM J.~Appl.~Dyn.~Syst.\ 14 (2015) 1454--1477.
\bibitem{Jef16} M.~R. Jeffrey, \emph{Exit from sliding in piecewise-smooth
flows: deterministic vs.\ determinacy-breaking}, Chaos 26 (2016) 033108.
\bibitem{BP21} J.~Bolte, E.~Pauwels, \emph{Conservative set-valued fields,
automatic differentiation, stochastic gradient methods and deep learning},
Math.~Program.\ 188 (2021).
\bibitem{APM20} A.~Aydinoglu, M.~Fazlyab, M.~Morari, M.~Posa, \emph{Stability
analysis of complementarity systems with neural network controllers},
arXiv:2011.07626.
\bibitem{Lu19} L.~Lu, Y.~Shin, Y.~Su, G.~E. Karniadakis, \emph{Dying ReLU and
initialization: theory and numerical examples}, arXiv:1903.06733.
\bibitem{ST98} J.~Sotomayor, M.~A. Teixeira, \emph{Regularization of
discontinuous vector fields}, in: Int.\ Conf.\ on Differential Equations
(Lisboa 1995), World Scientific, 1998, pp.~207--223.
\bibitem{LST07} J.~Llibre, P.~R. da Silva, M.~A. Teixeira, \emph{Regularization
of discontinuous vector fields on $\mathbb{R}^3$ via singular perturbation},
J.~Dynam.~Differential Equations 19 (2007) 309--331.
\bibitem{SDR23} M.~Sensi, M.~Desroches, S.~Rodrigues, \emph{Slow-fast dynamics
in a neurotransmitter release model: delayed response to a time-dependent input
signal}, Physica D (2023); arXiv:2302.12852.
\end{thebibliography}
\end{document}
